\section{Hilbert-valued row--log compression without orbit collapse}
\label{sec:tpc52-fourier-compressibility}

The preceding diagonal argument needs only the distinguished physical
direction.  It is nevertheless useful to locate that direction inside
the finite row--log banks of TPC--51.  The point of this section is
that smoothness gives a uniform finite approximation in the two
\emph{row--log} variables while the complete orbit and residual fiber
is left untouched.  In particular, none of the statements below
replaces the growing orbit space by a scalar space of polylogarithmic
rank.

\subsection{Tensor cosine projection}

Put \(T=\log 2\).  On \([0,T]\), let
\begin{equation}
 \chi_0(x):=T^{-1/2},
 \qquad
 \chi_k(x):=(2/T)^{1/2}\cos(\pi kx/T)\quad(k\geq1),
 \label{eq:tpc52-cosine-basis-one-dimensional}
\end{equation}
and set
\begin{equation}
 \chi_{k_1,k_2}(x,y):=\chi_{k_1}(x)\chi_{k_2}(y).
 \label{eq:tpc52-tensor-cosine-basis}
\end{equation}
This is an orthonormal basis of \(L^2([0,T]^2)\).  If \(\cK\) is a
Hilbert space and \(H:[0,T]^2\to\cK\) is Bochner measurable, define
\begin{equation}
 \Pi_{K_{\rm F}}^{\rm row}H
 :=\sum_{0\leq k_1,k_2\leq K_{\rm F}}
   \left(\int_{[0,T]^2}H(x,y)\chi_{k_1,k_2}(x,y)\,dx\,dy\right)
   \chi_{k_1,k_2}.
 \label{eq:tpc52-row-cosine-projection}
\end{equation}
The coefficient in parentheses is a vector in \(\cK\).

For later pointwise sampling, use the mixed regularity norm
\begin{equation}
 \|H\|_{\mathscr C^{2,2}}
 :=\sum_{0\leq a,b\leq2}
   \|\partial_x^a\partial_y^bH\|_{L^\infty([0,T]^2;\cK)}.
 \label{eq:tpc52-mixed-regularity-norm}
\end{equation}
The mixed derivatives in this definition are deliberately stronger
than the \(H^1\) regularity needed for the continuum \(L^2\) tail.  They
give a uniform bound at the arithmetic sampling nodes, including the
boundary of the dyadic log box.

\begin{theorem}[Hilbert-valued tensor cosine compression]
\label{thm:tpc52-Hilbert-cosine-compression}
There is a constant \(C_T<\infty\), independent of \(\cK\), such that
for every \(H\in C^{2,2}([0,T]^2;\cK)\) and \(K_{\rm F}\geq1\),
\begin{align}
 \|H-\Pi_{K_{\rm F}}^{\rm row}H\|_{L^2(e^{x+y}dxdy;\cK)}^2
 &\leq \frac{C_T}{(K_{\rm F}+1)^2}
       \|H\|_{\mathscr C^{2,2}}^2,
 \label{eq:tpc52-continuum-cosine-tail}\\
 \|H-\Pi_{K_{\rm F}}^{\rm row}H\|_{L^\infty([0,T]^2;\cK)}^2
 &\leq \frac{C_T}{(K_{\rm F}+1)^2}
       \|H\|_{\mathscr C^{2,2}}^2.
 \label{eq:tpc52-uniform-cosine-tail}
\end{align}
The same estimates hold for a finite or countable family
\((H_\alpha)_\alpha\) after summing both sides in \(\alpha\), whenever
the right side is finite.
\end{theorem}

\begin{proof}
For the \(L^2\) assertion, first use Lebesgue measure.  The tensor
cosines diagonalize the Neumann Laplacian.  Every omitted index has
  \(k_1>K_{\rm F}\) or \(k_2>K_{\rm F}\), and hence eigenvalue at least
  \((\pi(K_{\rm F}+1)/T)^2\).  Parseval, applied after scalar pairing with an
orthonormal basis of \(\cK\), gives
\[
 \|H-\Pi_{K_{\rm F}}^{\rm row}H\|_{L^2}^2
 \leq \frac{T^2}{\pi^2(K_{\rm F}+1)^2}
 \bigl(\|\partial_xH\|_{L^2}^2+
       \|\partial_yH\|_{L^2}^2\bigr).
\]
The density \(e^{x+y}\) is bounded above and below on the fixed box,
so this proves \eqref{eq:tpc52-continuum-cosine-tail}.

For the pointwise assertion, integrate the cosine coefficients twice
by parts in every nonzero coordinate.  The sine boundary terms vanish
at the first integration; the second integration produces only first
derivatives at the endpoints and integrals of second derivatives.
Consequently the Hilbert-valued coefficient \(\widehat H(k_1,k_2)\)
satisfies
\[
 \|\widehat H(k_1,k_2)\|_{\cK}
 \leq
 \frac{C_T\|H\|_{\mathscr C^{2,2}}}
      {(1+k_1)^2(1+k_2)^2}.
\]
The sum of this majorant over the complement of
\(\{0,\ldots,K_{\rm F}\}^2\) is
\(O_T((K_{\rm F}+1)^{-1})\).  Absolute convergence of
the cosine series, followed by squaring, proves
\eqref{eq:tpc52-uniform-cosine-tail}.  Every step is stable under
orthogonal summation in \(\cK\), which also proves the family version.
\end{proof}

\subsection{Arithmetic sampling and literal-mask deletion}

Let \(\cR_X\) be the complete prime--squarefree row universe and let
\(\cR_X(\cA)\subset\cR_X\) be one fixed positive-density admissible
row cell,
with logarithmic nodes
\[
 z_\gamma=(x_\gamma,y_\gamma)
 =\bigl(\log(\ell_\gamma/L),\log(d_\gamma/D)\bigr),
\]
 base weights \(w_\gamma=(\log\ell_\gamma)^{2e}\), where
 \(e\in\{0,1\}\), and total base energy
\(E_X=\sum_{\gamma\in\cR_X(\cA)}w_\gamma\).  We retain the inherited
base flatness
\begin{equation}
 E_X\asymp_{\cA}|\cR_X(\cA)|(\log L)^{2e},
 \qquad
 \max_{\gamma\in\cR_X(\cA)}w_\gamma\ll(\log L)^{2e}.
 \label{eq:tpc52-cosine-base-flatness}
\end{equation}
Here \(e=0\) is the unweighted cell and \(e=1\) is the
logarithmically weighted cell.  Let
\begin{equation}
 d\nu_{L,e}^{\rm ref}(x,y)
 =\frac{e^x(\log L+x)^{2e-1}}{Z_{L,e}}\,dx\,e^y\,dy
 \label{eq:tpc52-finite-row-reference-measure}
\end{equation}
be the finite-\(L\) reference probability measure of TPC--51.  For a
family \(H=(H_\alpha)_\alpha\), with values in an arbitrary Hilbert
space \(\cK_X\), define the normalized forms
\begin{align}
 \mathfrak q_X^0(H)
 &:=\frac1{E_X}\sum_\alpha
   \sum_{\gamma\in\cR_X(\cA)}
   w_\gamma\|H_\alpha(z_\gamma)\|_{\cK_X}^2,
 \label{eq:tpc52-unmasked-row-form}\\
 \mathfrak q_X^M(H)
 &:=\frac1{E_X}\sum_\alpha
   \sum_{\gamma\in\cR_X(\cA)}
   M_X(\alpha,\gamma)w_\gamma
   \|H_\alpha(z_\gamma)\|_{\cK_X}^2,
 \label{eq:tpc52-masked-row-form}\\
 \mathfrak q_{L,e}^{\rm ref}(H)
 &:=\sum_\alpha\int_{[0,T]^2}
   \|H_\alpha(x,y)\|_{\cK_X}^2
   \,d\nu_{L,e}^{\rm ref}(x,y).
 \label{eq:tpc52-reference-row-form}
\end{align}
Assume the literal mask has complete-universe row defect at most
\(\epsilon_X\).  Since the fixed cell has positive density, this implies
\begin{equation}
 \sup_\alpha
 \#\{\gamma\in\cR_X(\cA):M_X(\alpha,\gamma)=0\}
 \ll_{\cA}\epsilon_X|\cR_X(\cA)|.
 \label{eq:tpc52-cosine-cell-mask-defect}
\end{equation}
Put
\begin{equation}
 \Delta_X^{\rm mom}
 :=e^{-c\sqrt{\log L}}+D^{-1/2+\delta},
 \qquad 0<\delta<\frac12,
 \label{eq:tpc52-arithmetic-moment-error}
\end{equation}
where \(c>0\) is the fixed-modulus prime-progression constant.
The quantitative input proved in TPC--51 is that, for every fixed
\(A_0>0\), there is a fixed \(C_1\) such that
\begin{align}
 &\left|\int \e^{i(\xi x+\eta y)}\,d\nu_X
       -\int \e^{i(\xi x+\eta y)}\,d\nu_{L,e}^{\rm ref}\right|
 \notag\\[-2pt]
  &\hspace{25mm}\le C_{A_0}(1+|\xi|+|\eta|)^{C_1}\Delta_X^{\rm mom}
 \label{eq:tpc52-inherited-Fourier-moment}
\end{align}
uniformly for \(|\xi|+|\eta|\le C_T(\log X)^{A_0}\)
\citep{WangTPC51}.  Here \(\nu_X\) is the normalized empirical measure
associated with \eqref{eq:tpc52-unmasked-row-form}.  This displayed
estimate, rather than weak convergence alone, is the arithmetic hypothesis
used for a growing bank.

The error below has three logically different terms.  The first is a
smooth approximation tail.  The second is the arithmetic moment error
of the retained bank.  The third is the energy removed by the literal
mask.

\begin{theorem}[Sampled row--log compression ledger]
\label{thm:tpc52-sampled-cosine-compression}
Fix \(A_0>0\).  There are constants \(C,C_0<\infty\), depending on
the fixed row cell and \(A_0\) but not on \(X,K_{\rm F},\cK_X\), such that
the following holds whenever
\(1\leq K_{\rm F}\leq(\log X)^{A_0}\).  Let
\((H_\alpha)_\alpha\) have finite source support and put
\begin{equation}
 \mathfrak N(H)^2
 :=\sum_\alpha\|H_\alpha\|_{\mathscr C^{2,2}}^2,
 \qquad
 P_{K_{\rm F}}H:=(\Pi_{K_{\rm F}}^{\rm row}H_\alpha)_\alpha.
 \label{eq:tpc52-family-regularity-energy}
\end{equation}
Then
\begin{align}
 \mathfrak q_X^M(H-P_{K_{\rm F}}H)
 &\leq\mathfrak q_X^0(H-P_{K_{\rm F}}H)
 \leq \frac{C}{(K_{\rm F}+1)^2}\mathfrak N(H)^2,
 \label{eq:tpc52-discrete-relative-tail}\\
 \left|
  \mathfrak q_X^M(P_{K_{\rm F}}H)-
  \mathfrak q_{L,e}^{\rm ref}(P_{K_{\rm F}}H)
 \right|
 &\leq C\left(
   \Delta_X^{\rm mom}(K_{\rm F}+1)^{C_0}
   +\epsilon_X(K_{\rm F}+1)^2\right)\mathfrak N(H)^2.
 \label{eq:tpc52-bank-calibration-error}
\end{align}
In particular, the total approximation-and-calibration ledger obeys
\begin{align}
 &\mathfrak q_X^M(H-P_{K_{\rm F}}H)
 +\left|
  \mathfrak q_X^M(P_{K_{\rm F}}H)-
  \mathfrak q_{L,e}^{\rm ref}(P_{K_{\rm F}}H)
 \right|
 \notag\\[-2pt]
 &\qquad\leq C\left(
  (K_{\rm F}+1)^{-2}
  +\Delta_X^{\rm mom}(K_{\rm F}+1)^{C_0}
  +\epsilon_X(K_{\rm F}+1)^2\right)\mathfrak N(H)^2.
 \label{eq:tpc52-total-discrete-compression-error}
\end{align}
The exponent \(C_0\) is fixed.  One may take \(C_0=C_1+2\), where
\(C_1\) is the exponent in the prime--squarefree Fourier-moment error.
\end{theorem}

\begin{proof}
The mask only deletes nonnegative terms, so
\(\mathfrak q_X^M(H-P_{K_{\rm F}}H)
\leq\mathfrak q_X^0(H-P_{K_{\rm F}}H)\).  The normalized
row weights have bounded total mass.  Apply
\eqref{eq:tpc52-uniform-cosine-tail} for each \(\alpha\), then sum.
This proves \eqref{eq:tpc52-discrete-relative-tail}.

Write
\[
 P_{K_{\rm F}}H_\alpha(x,y)
 =\sum_{0\leq k_1,k_2\leq K_{\rm F}}
  a_{\alpha,k_1,k_2}\chi_{k_1,k_2}(x,y),
 \qquad a_{\alpha,k_1,k_2}\in\cK_X.
\]
Products of two tensor cosines are sums of \(O(1)\) exponentials with
frequencies \(O_T(K_{\rm F})\).  The fixed-modulus prime--squarefree moment
estimate therefore makes every sampled Gram entry differ from its
finite-\(L\) reference entry by at most
\[
 C(1+K_{\rm F})^{C_1}\Delta_X^{\rm mom}.
\]
There are \((K_{\rm F}+1)^2\) row--log modes.  The Schur bound gives the
operator error
\(C\Delta_X^{\rm mom}(K_{\rm F}+1)^{C_1+2}\).  This argument is unchanged
for Hilbert-valued coefficients because the scalar Gram is tensored
with \(I_{\cK_X}\).

For the mask term, the tensor-cosine Nikolskii estimate gives
\[
 \|P_{K_{\rm F}}H_\alpha\|_\infty^2
 \leq C_T(K_{\rm F}+1)^2
       \sum_{k_1,k_2}\|a_{\alpha,k_1,k_2}\|_{\cK_X}^2.
\]
At most an \(\epsilon_X\)-fraction of the complete opposite-row universe
is removed for each \(\alpha\).  The base-weight flatness on the fixed
row cell therefore bounds the deleted bank energy by
\(C\epsilon_X(K_{\rm F}+1)^2\sum_k\|a_{\alpha,k}\|^2\).  Sum in \(\alpha\).
Finally, Parseval and the definition of
\(\mathfrak N(H)\) bound the bank coefficient energy by
\(C\mathfrak N(H)^2\).  Combining the arithmetic and deletion errors
proves \eqref{eq:tpc52-bank-calibration-error} and hence
\eqref{eq:tpc52-total-discrete-compression-error}.
\end{proof}

\begin{proposition}[The uniformly active actual field meets the normalization]
\label{prop:tpc52-actual-field-normalization}
Assume the compact \(C^\infty\) profile family of
\eqref{eq:tpc52-profile-factorization}, the carrier bounds
\eqref{eq:tpc52-carrier-normalization}, and
Assumptions~\ref{ass:tpc52-profile-nondegeneracy} and
\ref{ass:tpc52-exact-profile-dictionary}.  Let \(\cU_X\) be a
finite set of source/cofactor pairs \((\alpha,r)\), along the inherited
regime \(J,L\to\infty\).  Let \(a_\alpha\) be the actual source coefficient,
independent of the retained cofactor label \(r\), and let the source
parameter be \(\beta_\alpha\in\mathfrak B\).  Define the retained field,
indexed by the
source label used in the literal mask, by
\begin{equation}
 H_{\alpha,X}(x,y)
 :=\bigoplus_{r:(\alpha,r)\in\cU_X}a_\alpha
   \bigoplus_{a\in\cA_0}
   \bigoplus_{j\in\cJ_{X,a}}
   F_{\beta_\alpha,a}(x,y,j/J)e_{r,a,j}.
 \label{eq:tpc52-actual-retained-field}
\end{equation}
The coordinate \(e_{r,a,j}\) may be tensored with any retained residual
basis vector; no coherent sum over those labels is taken.  Put
\begin{equation}
 \mathcal C_X^{\rm act}
 :=J\sum_{(\alpha,r)\in\cU_X}|a_\alpha|^2
       \mathscr B(\beta_\alpha).
 \label{eq:tpc52-actual-field-carrier-normalization}
\end{equation}
Then, uniformly in the coefficients and parameters,
\begin{align}
 \mathfrak N(H_X)^2
 &\le C_{\rm reg}\mathcal C_X^{\rm act},
 \label{eq:tpc52-actual-field-regularity}\\
 \mathfrak q_{L,e}^{\rm ref}(H_X)
 &\ge \left(\rho_*+o(1)\right)\mathcal C_X^{\rm act}.
 \label{eq:tpc52-actual-field-reference-overlap}
\end{align}
Consequently, if \(K_{{\rm F},X}\to\infty\), then for all sufficiently large
\(X\),
\begin{equation}
 \mathfrak q_{L,e}^{\rm ref}(P_{K_{{\rm F},X}}H_X)
 \ge c_{\rm ov}\mathcal C_X^{\rm act}
 \label{eq:tpc52-actual-projected-reference-overlap}
\end{equation}
with a fixed \(c_{\rm ov}>0\).  Thus the two normalization hypotheses in
\eqref{eq:tpc52-relative-tail-normalization} hold for the actual retained
field, not merely for a model family.
\end{proposition}

\begin{proof}
For every derivative occurring in
\eqref{eq:tpc52-mixed-regularity-norm}, orthogonality of the orbit
coordinates gives
\[
 \sup_{x,y}\|\partial_x^v\partial_y^wH_{\alpha,X}(x,y)\|^2
 \le C J\sum_{r:(\alpha,r)\in\cU_X}|a_\alpha|^2.
\]
There are only finitely many derivative orders and direct cells.  Sum in
\(\alpha\), use Cauchy--Schwarz over those derivative orders, and then use
\(\mathscr B\ge b_-\) to obtain
\eqref{eq:tpc52-actual-field-regularity}.

The progression Riemann sum in \(j/J\), uniformly for the compact profile
family, and uniform convergence of the finite-\(L\) reference density to
\(e^{x+y}\) give
\[
 \mathfrak q_{L,e}^{\rm ref}(H_X)
 =J\sum_{(\alpha,r)\in\cU_X}|a_\alpha|^2
   \bigl(\mathscr P(\beta_\alpha)+o(1)\bigr).
\]
The error is uniform in \(\alpha\) and \(\beta_\alpha\); after weighting and summing,
it is \(o(\mathcal C_X^{\rm act})\) because
\(b_-\le\mathscr B\le b_+\).  Apply
\(\mathscr P(\beta_\alpha)\ge\rho_*\mathscr B(\beta_\alpha)\) to prove
\eqref{eq:tpc52-actual-field-reference-overlap}.

Finally, the weighted continuum estimate
\eqref{eq:tpc52-continuum-cosine-tail} and
\eqref{eq:tpc52-actual-field-regularity} imply
\[
 \|H_X-P_{K_{{\rm F},X}}H_X\|_{L^2(d\nu_{L,e}^{\rm ref};\cK_X)}
 \le C(K_{{\rm F},X}+1)^{-1}(\mathcal C_X^{\rm act})^{1/2}.
\]
The triangle inequality, followed by \(K_{{\rm F},X}\to\infty\), gives
\eqref{eq:tpc52-actual-projected-reference-overlap}.
\end{proof}

\begin{corollary}[Polylogarithmic relative tail]
\label{cor:tpc52-polylog-relative-tail}
Suppose a retained actual family satisfies, for one declared
normalization \(\mathcal C_X>0\),
\begin{equation}
 \mathfrak N(H_X)^2\leq C_{\rm reg}\mathcal C_X,
 \qquad
 \mathfrak q_{L,e}^{\rm ref}(P_{K_{\rm F}}H_X)\geq c_{\rm ov}\mathcal C_X
 \label{eq:tpc52-relative-tail-normalization}
\end{equation}
with fixed \(C_{\rm reg}<\infty\) and \(c_{\rm ov}>0\).  Let
\(K_{\rm F}=\lfloor(\log X)^A\rfloor\), where \(A>0\) is fixed, and apply
\cref{thm:tpc52-sampled-cosine-compression} with a fixed
\(A_0\ge A\).  If
\begin{equation}
 \Delta_X^{\rm mom}K_{\rm F}^{C_0}=o(1),
 \qquad
 \epsilon_X(K_{\rm F}+1)^2=o(1),
 \label{eq:tpc52-polylog-conditions}
\end{equation}
then \(P_{K_{\rm F}}H_X\) has masked energy comparable to
\(\mathcal C_X\), while
\begin{equation}
 \mathfrak q_X^M(H_X-P_{K_{\rm F}}H_X)
 \leq\mathfrak q_X^0(H_X-P_{K_{\rm F}}H_X)
 \ll (\log X)^{-2A}\mathcal C_X=o(\mathcal C_X).
 \label{eq:tpc52-polylog-tail-conclusion}
\end{equation}
Thus the retained bank verifies the relative-tail alternative for this
actual family.

If \(L,D\geq X^{c_*}\) for some fixed \(c_*>0\),
\(0<\delta<1/2\), and \(\epsilon_X\leq X^{-\eta+o(1)}\) for some fixed
\(\eta>0\), conditions \eqref{eq:tpc52-polylog-conditions} hold for
every fixed \(A>0\).
\end{corollary}

\begin{proof}
Insert \(K_{\rm F}=(\log X)^A\) in
\eqref{eq:tpc52-total-discrete-compression-error}.  The exponential
prime error absorbs every fixed log power, the squarefree term has a
fixed power saving because \(\delta<1/2\), and the mask term has a
fixed power saving by hypothesis.  The lower bound in
\eqref{eq:tpc52-relative-tail-normalization} absorbs the resulting
\(o(\mathcal C_X)\) calibration error.  Equation
\eqref{eq:tpc52-polylog-tail-conclusion} is
\eqref{eq:tpc52-discrete-relative-tail}.
\end{proof}

By \cref{prop:tpc52-actual-field-normalization}, this corollary applies to
the uniformly active provenance-fixed field with
\(\mathcal C_X=\mathcal C_X^{\rm act}\).  Hence the relative-tail conclusion
used in the introduction is an actual retained-label statement under the
same exact-dictionary and profile-nondegeneracy hypotheses as the diagonal
bridge.

\subsection{What is and is not compressed}

For the actual coefficient field, take
\begin{equation}
 \cK_X
 =\bigoplus_{r,j}\cK_{X,r,j},
 \label{eq:tpc52-retained-orbit-fiber}
\end{equation}
where the summands retain the cofactor, orbit, and residual data required
by \eqref{eq:tpc52-coefficient-space}; the source remains the family index,
while the opposite-row coordinates \((x,y)\) are the arguments of the field.  The
operator used above is exactly
\begin{equation}
 \Pi_{K_{\rm F}}^{\rm row}\otimes I_{\cK_X}.
 \label{eq:tpc52-row-only-projection}
\end{equation}
It has \((K_{\rm F}+1)^2\) row--log modes, but every coefficient of those modes
is a full vector in \(\cK_X\).  In particular,
\begin{equation}
 \dim\operatorname{ran}(\Pi_{K_{\rm F}}^{\rm row}\otimes I_{\cK_X})
 =(K_{\rm F}+1)^2\dim\cK_X
 \label{eq:tpc52-no-orbit-rank-compression}
\end{equation}
when \(\cK_X\) is finite-dimensional.  If the physical orbit rank is of
order \(J\), that factor is still present.  Hence
\cref{thm:tpc52-sampled-cosine-compression} is compatible with scalar
translate saturation: it compresses smooth dependence on the two row
logs and makes no fixed-rank assertion about the orbit direction.
